%! TeX program = xelatex
\documentclass[12pt,national]{teaching}

\renewcommand\thesection{Exercice~\arabic{section}~:}
\renewcommand\thesubsection{Partie~\Roman{subsection}~:}

\newlength\cwi
\newlength\cwii
\newlength\cwiii

\def\b{16cm}
\def\c{2cm}
\def\a{\linewidth-\b-\c}

\begin{document}

\section{Chimie (7 points)}
\subsection{Nickelage d'une lame d'aluminium}

\setlength\cwii{\b}
\setlength\cwiii{\c}
\setlength\cwi{\a}
\begin{tcbitemize}[raster columns=3,raster equal height=rows,
        raster force size=false,
        raster equal skip=-1.0pt,
        raster column 1/.style={width=\cwi,
            fontupper=\bfseries},
        raster column 2/.style={width=\cwii},
        raster column 3/.style={width=\cwiii,fontupper=\bfseries,
            halign=center,valign=center},
        ]
    %----------------------------------------------------------------
    \tcbitem 1-
    \tcbitem $\ce{Ni^{2+}_{(aq)} + 2e- -> Ni_{(s)}}$
    \tcbitem (0,5pt)
    %----------------------------------------------------------------
    \tcbitem[raster multirow=2] 2-
    \tcbitem[raster multicolumn=2,raster multirow=2,blankest]
        \begin{tcbitemize}[raster rows=2,raster columns=2,
               raster height=\tcbtextheight,
               raster column 1/.style={width=\cwii},
               raster column 2/.style={width=\cwiii,fontupper=\bfseries}
            ]
            \tcbitem la réduction a lieu à la cathode.
            \tcbitem (0,25pt)
            %--------------------------------------------------------
            \tcbitem la lame d'aluminium est la cathode
            \tcbitem (0,25pt)
        \end{tcbitemize}
    %----------------------------------------------------------------
    \tcbitem[raster multirow=2] 3-
    \tcbitem[raster multicolumn=2,raster multirow=2,blankest]
        \begin{tcbitemize}[raster rows=2,raster columns=2,
               raster height=\tcbtextheight,
               raster column 1/.style={width=\cwii},
               raster column 2/.style={width=\cwiii,fontupper=\bfseries}
            ]
            \tcbitem $Q=I\Delta t=nF$ \quad avec $n=2x$ \quad
                     donc $x=\frac{I\Delta t}{2F}$
                     
                     $n(\ce{Ni})=x=\frac{I\Delta t}{2F}$

                     $m(\ce{Ni})=n(\ce{Ni})M(\ce{Ni})$ \quad donc
                     $m(\ce{Ni})=\frac{I\Delta t\,M(\ce{Ni})}{2F}$
            \tcbitem (0,5pt)
            %--------------------------------------------------------
            \tcbitem $m(\ce{Ni})=
                     \frac{\num{100e-3}\times 25\times 60\times 58,7}
                     {2\times\num{9.65e4}}=\qty{4,56e-2}{\g}=\qty{45,6}{\mg}$
            \tcbitem (0,25pt)
        \end{tcbitemize}
    %----------------------------------------------------------------
    \tcbitem[raster multirow=2] 4-
    \tcbitem[raster multicolumn=2,raster multirow=2,blankest]
        \begin{tcbitemize}[raster rows=2,raster columns=2,
               raster height=\tcbtextheight,
               raster column 1/.style={width=\cwii},
               raster column 2/.style={width=\cwiii,fontupper=\bfseries}
            ]
            \tcbitem $\rho=\frac{m}{V}=\frac{m}{s\times e}$ \quad
                     donc\, $e=\frac{m}{\rho s}$
            \tcbitem (0,25pt)
            %--------------------------------------------------------
            \tcbitem $e=\frac{\num{4.56e-2}}{8,9\times 50}=
                     \qty{1.025e-4}{\cm}=\qty{1.025e-6}{\m}\simeq\qty{1.0}{\um}$
            \tcbitem (0,25pt)
        \end{tcbitemize}
\end{tcbitemize}


\subsection{Étude de quelques réactions d'un acide AH}


\setlength\cwii{\b}
\setlength\cwiii{\c}
\setlength\cwi{\a}
\begin{tcbitemize}[raster columns=3,raster equal height=rows,
        raster force size=false,
        raster equal skip=-1.0pt,
        raster column 1/.style={width=\cwi,
            fontupper=\bfseries},
        raster column 2/.style={width=\cwii},
        raster column 3/.style={width=\cwiii,fontupper=\bfseries,
            halign=center,valign=center},
        ]
    %----------------------------------------------------------------
    \tcbitem 1-
    \tcbitem $\ce{AH_{(aq)} + H2O_{(\ell)} <=> A^-_{(aq)} + H3O^+_{(aq)}}$
    \tcbitem (0,5pt)
    %----------------------------------------------------------------
    \tcbitem 2-
    \tcbitem $\tau=\frac{x_{\eq}}{x_{\max}}=
             \frac{[\ce{H3O+}]_{\eq}\cdot V}{C_A\cdot V}=
             \frac{[\ce{H3O+}]_{\eq}}{C_A}=
             \frac{10^{-\mathrm{pH}}}{C_A}$
    \tcbitem (0,5pt)
    %----------------------------------------------------------------
    \tcbitem 3-1-
    \tcbitem d'après la courbe, pour $\mathrm{pH}=2,9$, on lit~: $-\log C_A=1$ \\
             donc\, $C_A=10^{-1}=\qty{0.10}{\M}$
    \tcbitem (0,5pt)
    %----------------------------------------------------------------
    \tcbitem[raster multirow=2] 3-2-
    \tcbitem[raster multicolumn=2,raster multirow=2,blankest]
        \begin{tcbitemize}[raster rows=2,raster columns=2,
               raster height=\tcbtextheight,
               raster column 1/.style={width=\cwii},
               raster column 2/.style={width=\cwiii,fontupper=\bfseries}
            ]
            \tcbitem $\tau_{1}=\frac{10^{-\mathrm{pH}_{1}}}{C_{A_1}}=
                     \frac{10^{-2,9}}{0,10}=\num{1.26e-2}=\qty{1.26}{\percent}$
            \tcbitem (0,25pt)
            %--------------------------------------------------------
            \tcbitem $\tau_{2}=\frac{10^{-\mathrm{pH}_{2}}}{C_{A_2}}=
                     \frac{10^{-3,15}}{10^{-1,5}}=
                     \num{2.24e-2}=\qty{2.24}{\percent}$
            \tcbitem (0,5pt)
        \end{tcbitemize}
    %----------------------------------------------------------------
    \tcbitem[raster multirow=2] 3-3-
    \tcbitem[raster multicolumn=2,raster multirow=2,blankest]
        \begin{tcbitemize}[raster rows=2,raster columns=2,
               raster height=\tcbtextheight,
               raster column 1/.style={width=\cwii},
               raster column 2/.style={width=\cwiii,fontupper=\bfseries}
            ]
            \tcbitem $C_{A_2}<C_{A_1} \quad \Rightarrow\, \tau_2>\tau_1$
            \tcbitem (0,25pt)
            %--------------------------------------------------------
            \tcbitem la dilution augmente le taux d'avancement final
            \tcbitem (0,25pt)
        \end{tcbitemize}
    %----------------------------------------------------------------
    \tcbitem[raster multirow=2] 4-1-
    \tcbitem[raster multicolumn=2,raster multirow=2,blankest]
        \begin{tcbitemize}[raster rows=2,raster columns=2,
               raster height=\tcbtextheight,
               raster column 1/.style={width=\cwii},
               raster column 2/.style={width=\cwiii,fontupper=\bfseries}
            ]
            \tcbitem $\mathrm{K_A}=
                     \frac{[\ce{A-}]_{\eq}\cdot [\ce{H3O+}]_{\eq}}
                     {[\ce{AH}]_{\eq}}\simeq
                     \frac{[\ce{H3O+}]_{\eq}^2}{C_A}$ \quad
                     $\Rightarrow\, [\ce{H3O+}]_{\eq}=\sqrt{\mathrm{K_A}C_A}$

                     $-\log[\ce{H3O+}]_{\eq}=
                     -\frac12\log\mathrm{K_A}-\frac12\log C_A$ \quad
                     $\Rightarrow\,\mathrm{pH}=
                     \frac12\mathrm{pK_A}-\frac12\log C_A$

                     D'où \, $\mathrm{pH}=a\log C_A+b$ \, avec \,
                     $a=-\frac12$ et $b=\frac12\mathrm{pK_A}$
            \tcbitem (0,75pt)
            %--------------------------------------------------------
            \tcbitem 
            \tcbitem (0,0pt)
        \end{tcbitemize}
    %----------------------------------------------------------------
    \tcbitem[raster multirow=2] 4-2-
    \tcbitem[raster multicolumn=2,raster multirow=2,blankest]
        \begin{tcbitemize}[raster rows=2,raster columns=2,
               raster height=\tcbtextheight,
               raster column 1/.style={width=\cwii},
               raster column 2/.style={width=\cwiii,fontupper=\bfseries}
            ]
            \tcbitem pour $(\mathrm{pH}~;~C_{A})=(2,9~;~\qty{0.10}{\M})$

                     $\mathrm{pK_A}=2\mathrm{pH}+\log C_A=
                     2\times 2,9+\log(0,10)=4,8$
            \tcbitem (0,5pt)
            %--------------------------------------------------------
            \tcbitem l'acide $\ce{AH}$ est $\ce{CH3COOH}$
            \tcbitem (0,25pt)
        \end{tcbitemize}
    %----------------------------------------------------------------
    \tcbitem[raster multirow=2] 4-3-
    \tcbitem[raster multicolumn=2,raster multirow=2,blankest]
        \begin{tcbitemize}[raster rows=2,raster columns=2,
               raster height=\tcbtextheight,
               raster column 1/.style={width=\cwii},
               raster column 2/.style={width=\cwiii,fontupper=\bfseries}
            ]
            \tcbitem $\mathrm{pH}=-\frac12\log C_A+\frac12\mathrm{pK_A}=
                     -\frac12\times\log(\num{1e-2})+\frac12\times 4,8=3,4$
            \tcbitem (0,25pt)
            %--------------------------------------------------------
            \tcbitem $\mathrm{pH}<\mathrm{pK_A}$, \;
                     alors l'acide $\ce{AH}$ prédomine
            \tcbitem (0,25pt)
        \end{tcbitemize}
\end{tcbitemize}

\section{Transformations nucléaires (2,5 points)}

\setlength\cwii{\b}
\setlength\cwiii{\c}
\setlength\cwi{\a}
\begin{tcbitemize}[raster columns=3,raster equal height=rows,
        raster force size=false,
        raster equal skip=-1.0pt,
        raster column 1/.style={width=\cwi,
            fontupper=\bfseries},
        raster column 2/.style={width=\cwii},
        raster column 3/.style={width=\cwiii,fontupper=\bfseries,
            halign=center,valign=center},
        ]
    %----------------------------------------------------------------
    \tcbitem[raster multirow=2] 1-
    \tcbitem[raster multicolumn=2,raster multirow=2,blankest]
        \begin{tcbitemize}[raster rows=2,raster columns=2,
               raster height=\tcbtextheight,
               raster column 1/.style={width=\cwii},
               raster column 2/.style={width=\cwiii,fontupper=\bfseries}
            ]
            \tcbitem $52$ protons 
            \tcbitem (0,25pt)
            %--------------------------------------------------------
            \tcbitem $138-52=86$ neutrons
            \tcbitem (0,25pt)
        \end{tcbitemize}
    %----------------------------------------------------------------
    \tcbitem[raster multirow=2] 2-
    \tcbitem[raster multicolumn=2,raster multirow=2,blankest]
        \begin{tcbitemize}[raster rows=2,raster columns=2,
               raster height=\tcbtextheight,
               raster column 1/.style={width=\cwii},
               raster column 2/.style={width=\cwiii,fontupper=\bfseries}
            ]
            \tcbitem conservation du nombre de masse~: $235+1=95+138+x$
            \tcbitem (0,25pt)
            %--------------------------------------------------------
            \tcbitem $\Rightarrow\, x=3$
            \tcbitem (0,25pt)
        \end{tcbitemize}
    %----------------------------------------------------------------
    \tcbitem 3-
    \tcbitem $E_{\ell}\left(\ce{^{138}_{52}Te}\right)=
             \left[52m_p+86m_n-m\left(\ce{^{138}_{52}Te}\right)\right]\times c^2$ 

             $E_{\ell}\left(\ce{^{138}_{52}Te}\right)=
             \left[52\times1,00727+86\times1,00866-137,92922\right]\times931,49=
             \qty{1111.8}{\MeV}$

    \tcbitem (0,5pt)
    %----------------------------------------------------------------
    \tcbitem 4-
    \tcbitem 
             \begin{itemize}
                 \item $\frac{E_{\ell}\left(\ce{^{138}_{52}Te}\right)}{138}=
                       \frac{1111.8}{138}=\qty{8.06}{\MeV\per\mathrm{nucléon}}$
                 \item $\frac{E_\ell\left(\ce{^{95}_{40}Zr}\right)}{95}=
                       \frac{800,214}{95}=\qty{8.42}{\MeV\per\mathrm{nucléon}}$
             \end{itemize}
             $\frac{E_\ell\left(\ce{^{95}_{40}Zr}\right)}{95}>
             \frac{E_{\ell}\left(\ce{^{138}_{52}Te}\right)}{138}$ \quad
             $\Rightarrow\,
             \ce{^{95}_{40}Zr}$ est plus stable que $\ce{^{138}_{52}Te}$
    \tcbitem (0,5pt)
    %----------------------------------------------------------------
    \tcbitem 5-
    \tcbitem $|\Delta E_{1}|=\left|E_\ell\left(\ce{^{235}_{92}U}\right)-
             E_\ell\left(\ce{^{95}_{40}Zr}\right)-
             E_\ell\left(\ce{^{138}_{52}Te}\right)\right|$

             $|\Delta E_{1}|=|1735,7-800,214-1111,8|=\qty{176.314}{\MeV}$

             $|\Delta E|=N|\Delta E_{1}|=\frac{mN_{A}}{M(\ce{U})}|\Delta E_{1}|$

             $|\Delta E|=\frac{\num{1e3}\times\num{6.02e23}}{235}\times176.314=
             \qty{4.5166e26}{\MeV}$

             $|\Delta E|=\num{4.5166e26}\times\num{1.60e-13}=
             \qty{7.2266e13}{\joule}$
    \tcbitem (0,5pt)
\end{tcbitemize}



\section{Électricité (5 points)}

{
\renewcommand\thesubsection{\arabic{subsection}-}
\subsection{Charge d'un condensateur}

\setlength\cwii{\b}
\setlength\cwiii{\c}
\setlength\cwi{\a}
\begin{tcbitemize}[raster columns=3,raster equal height=rows,
        raster force size=false,
        raster equal skip=-1.0pt,
        raster column 1/.style={width=\cwi,
            fontupper=\bfseries},
        raster column 2/.style={width=\cwii},
        raster column 3/.style={width=\cwiii,fontupper=\bfseries,
            halign=center,valign=center},
        ]
    %----------------------------------------------------------------
    \tcbitem 1-1-
    \tcbitem loi des mailles~: $E=u_R+u_C$ \\
             or~: $u_R=R_0i$ et $i=C_0\frac{\mathrm{d}u_C}{\mathrm{d}t}$ \\
             donc~: $R_0C_0\frac{\mathrm{d}u_C}{\mathrm{d}t}+u_C=E$
    \tcbitem (0,5pt)
    %----------------------------------------------------------------
    \tcbitem[raster multirow=2] 1-2-
    \tcbitem[raster multicolumn=2,raster multirow=2,blankest]
        \begin{tcbitemize}[raster rows=2,raster columns=2,
               raster height=\tcbtextheight,
               raster column 1/.style={width=\cwii},
               raster column 2/.style={width=\cwiii,fontupper=\bfseries}
            ]
            \tcbitem $u_C(t)=E\left(1-e^{-\frac{t}{\tau}}\right)$ \, et \,
                     $\frac{\mathrm{d}u_C}{\mathrm{d}t}=
                     \frac{E}{\tau}e^{-\frac{t}{\tau}}$

                     $\Rightarrow\,
                     R_0C_0\times\frac{E}{\tau}e^{-\frac{t}{\tau}}+
                     E-E\times e^{-\frac{t}{\tau}}=E$ \\
                     $\Rightarrow\,
                     Ee^{-\frac{t}{\tau}}\left(\frac{R_0C_0}{\tau}-1\right)=0$\\
                     pour tout $t$,
                     $\frac{R_0C_0}{\tau}-1=0$
            \tcbitem (0,25pt)
            %--------------------------------------------------------
            \tcbitem donc $\tau=R_0C_0$
            \tcbitem (0,25pt)
        \end{tcbitemize}
    %----------------------------------------------------------------
    \tcbitem 1-3-
    \tcbitem $u_R(t)=E-u_C(t)=E-E\left(1-e^{-\frac{t}{\tau}}\right)=
             E e^{-\frac{t}{\tau}}$
    \tcbitem (0,5pt)
    %----------------------------------------------------------------
    \tcbitem 1-4-
    \tcbitem $\ln(u_R)=\ln(E)-\frac{t}{\tau}$
    \tcbitem (0,25pt)
    %----------------------------------------------------------------
    \tcbitem[raster multirow=2] 1-5-
    \tcbitem[raster multicolumn=2,raster multirow=2,blankest]
        \begin{tcbitemize}[raster rows=2,raster columns=2,
               raster height=\tcbtextheight,
               raster column 1/.style={width=\cwii},
               raster column 2/.style={width=\cwiii,fontupper=\bfseries}
            ]
            \tcbitem $\ln(E)=2,2$ \quad
                     $\Rightarrow\, E=e^{2,2}=\qty{9.0}{\volt}$
            \tcbitem (0,25pt)
            %--------------------------------------------------------
            \tcbitem $-\frac{1}{\tau}=\frac{2,2-0}{0-\num{55e-5}}=
                     -\qty{4000}{\per\s} \quad\Rightarrow\,
                     \tau=\qty{2.5e-4}{\s}$

                     $\tau=R_0C_0$ \quad
                     $\Rightarrow\, C_0=\frac{\tau}{R_0}=
                     \frac{\num{2.5e-4}}{\num{100e3}}=
                     \qty{2.5e-9}{\F}=\qty{2.5}{\nF}$
            \tcbitem (0,5pt)
        \end{tcbitemize}
    %----------------------------------------------------------------
    \tcbitem[raster multirow=2] 1-6-
    \tcbitem[raster multicolumn=2,raster multirow=2,blankest]
        \begin{tcbitemize}[raster rows=2,raster columns=2,
               raster height=\tcbtextheight,
               raster column 1/.style={width=\cwii},
               raster column 2/.style={width=\cwiii,fontupper=\bfseries}
            ]
            \tcbitem à $t=\qty{55e-5}{\s}$, on lit~: $\ln(u_R)=0$ \,
                     donc $u_R=1$ \, ainsi $u_C=E-u_R=9-1=\qty{8}{\V}$

                     $E_e=\frac12 C_0u_C^2=\frac12\times\num{2.5e-9}\times 8^2$
            \tcbitem (0,25pt)
            %--------------------------------------------------------
            \tcbitem $E_e=\qty{8.0e-8}{\joule}$
            \tcbitem (0,25pt)
        \end{tcbitemize}
\end{tcbitemize}

\subsection{Réception d'un signal modulé en amplitude}

\setlength\cwii{\b}
\setlength\cwiii{\c}
\setlength\cwi{\a}
\begin{tcbitemize}[raster columns=3,raster equal height=rows,
        raster force size=false,
        raster equal skip=-1.0pt,
        raster column 1/.style={width=\cwi,
            fontupper=\bfseries},
        raster column 2/.style={width=\cwii},
        raster column 3/.style={width=\cwiii,fontupper=\bfseries,
            halign=center,valign=center},
        ]
    %----------------------------------------------------------------
    \tcbitem 2-1-
    \tcbitem C
    \tcbitem (0,5pt)
    %----------------------------------------------------------------
    \tcbitem[raster multirow=2] 2-2-
    \tcbitem[raster multicolumn=2,raster multirow=2,blankest]
        \begin{tcbitemize}[raster rows=2,raster columns=2,
               raster height=\tcbtextheight,
               raster column 1/.style={width=\cwii},
               raster column 2/.style={width=\cwiii,fontupper=\bfseries}
            ]
            \tcbitem $F=\frac{1}{2\pi\sqrt{L_0C_0}}$
            \tcbitem (0,5pt)
            %--------------------------------------------------------
            \tcbitem $F=\frac{1}{2\pi\sqrt{\num{1e-3}\times\num{2.5e-9}}}
                     \simeq \qty{1.0e5}{\hertz}=\qty{100}{\kilo\hertz}$
            \tcbitem (0,25pt)
        \end{tcbitemize}
    %----------------------------------------------------------------
    \tcbitem[raster multirow=2] 2-3-
    \tcbitem[raster multicolumn=2,raster multirow=2,blankest]
        \begin{tcbitemize}[raster rows=2,raster columns=2,
               raster height=\tcbtextheight,
               raster column 1/.style={width=\cwii},
               raster column 2/.style={width=\cwiii,fontupper=\bfseries}
            ]
            \tcbitem pour une bonne détection d'enveloppe, il faut~:
                     $T_p \ll RC < T_m$ \\
                     avec~: $T_p=\frac{1}{F}=\qty{1e-5}{\s}=\qty{0.01}{\ms}$ \\
                     d'après la figure 4~: $T_m=2\times 0,4=\qty{0.8}{\ms}$ \\
                     on teste les condensateurs~:
                     \begin{itemize}
                         \item $C=\qty{120}{\nF}
                               \quad\Rightarrow\, RC=\qty{1.2}{\ms}$
                         \item $C=\qty{100}{\nF}
                               \quad\Rightarrow\, RC=\qty{1.0}{\ms}$
                         \item $C=\qty{50}{\nF}
                               \quad\Rightarrow\, RC=\qty{0.5}{\ms}$
                     \end{itemize}
            \tcbitem (0,5pt)
            %--------------------------------------------------------
            \tcbitem $C=\qty{50}{\nano\farad}$
            \tcbitem (0,25pt)
        \end{tcbitemize}
\end{tcbitemize}
}

\section{Mécanique (5,5 points)}

\subsection{Chute verticale d'une balle}

\setlength\cwii{\b}
\setlength\cwiii{\c}
\setlength\cwi{\a}
\begin{tcbitemize}[raster columns=3,raster equal height=rows,
        raster force size=false,
        raster equal skip=-1.0pt,
        raster column 1/.style={width=\cwi,
            fontupper=\bfseries},
        raster column 2/.style={width=\cwii},
        raster column 3/.style={width=\cwiii,fontupper=\bfseries,
            halign=center,valign=center},
        ]
    %----------------------------------------------------------------
    \tcbitem 1-
    \tcbitem d'après la deuxième loi de Newton~: $\sum \vec F = m\vec a$ \\
             on projette sur $Oz$~: $P-f=ma_z$ \\
             or \, $P=mg,\qquad f=\mu v^{2}=\mu v_z^2,\qquad
             a_z=\frac{\mathrm{d}v_z}{\mathrm{d}t}$ \\
             donc \, $mg-\mu v_z^2=m\frac{dv_z}{dt}$ \quad
             $\Rightarrow\,
             \frac{\mathrm{d}v_z}{\mathrm{d}t}+\frac{\mu}{m}v_z^2=g$
    \tcbitem (0,5pt)
    %----------------------------------------------------------------
    \tcbitem[raster multirow=2] 2-
    \tcbitem[raster multicolumn=2,raster multirow=2,blankest]
        \begin{tcbitemize}[raster rows=2,raster columns=2,
               raster height=\tcbtextheight,
               raster column 1/.style={width=\cwii},
               raster column 2/.style={width=\cwiii,fontupper=\bfseries}
            ]
            \tcbitem à la vitesse limite~: $\frac{dv_z}{dt}=0$
                    donc $\frac{\mu}{m}v_{z\lim}^2=g$ \quad
                    $\Rightarrow\, v_{z\lim}=\sqrt{\frac{mg}{\mu}}$
            \tcbitem (0,5pt)
            %--------------------------------------------------------
            \tcbitem $\mu=\frac{mg}{v_{z\lim}^2}=
                     \frac{\num{2.45e-3}\times 10}{7,00^2}=
                     \qty{5.0e-4}{\kg\per\m}$
            \tcbitem (0,25pt)
        \end{tcbitemize}
    %----------------------------------------------------------------
    \tcbitem[raster multirow=2] 3-
    \tcbitem[raster multicolumn=2,raster multirow=2,blankest]
        \begin{tcbitemize}[raster rows=2,raster columns=2,
               raster height=\tcbtextheight,
               raster column 1/.style={width=\cwii},
               raster column 2/.style={width=\cwiii,fontupper=\bfseries}
            ]
            \tcbitem $\frac{v_z(t_{2})-v_z(t_{1})}{t_{2}-t_{1}}+
                     \frac{\mu}{m}v_z^2(t_{1})=g$ \quad
                     $\Rightarrow\, v_z(t_2)=v_z(t_1)+
                     0,07\left(g-\frac{\mu}{m}v_z^2(t_1)\right)$
            \tcbitem (0,5pt)
            %--------------------------------------------------------
            \tcbitem $v_z(t_2)=5,50+
                     0,07\times\left(10-\frac{\num{5.0e-4}}{\num{2.45e-3}}\times
                     5,50^2\right)\simeq\qty{5.77}{\v}$
            \tcbitem (0,25pt)
        \end{tcbitemize}
\end{tcbitemize}

\subsection{Mouvement d'un système oscillant}

\setlength\cwii{\b}
\setlength\cwiii{\c}
\setlength\cwi{\a}
\begin{tcbitemize}[raster columns=3,raster equal height=rows,
        raster force size=false,
        raster equal skip=-1.0pt,
        raster column 1/.style={width=\cwi,
            fontupper=\bfseries},
        raster column 2/.style={width=\cwii},
        raster column 3/.style={width=\cwiii,fontupper=\bfseries,
            halign=center,valign=center},
        ]
    %----------------------------------------------------------------
    \tcbitem 1-
    \tcbitem $E_{pe}=\frac12 kx^2$
    \tcbitem (0,5pt)
    %----------------------------------------------------------------
    \tcbitem[raster multirow=2] 2-
    \tcbitem[raster multicolumn=2,raster multirow=2,blankest]
        \begin{tcbitemize}[raster rows=2,raster columns=2,
               raster height=\tcbtextheight,
               raster column 1/.style={width=\cwii},
               raster column 2/.style={width=\cwiii,fontupper=\bfseries}
            ]
            \tcbitem $\frac12k=\frac{(4-2)\times10^{-3}}{(2-1)\times 10^{-4}}=
                     \qty{20}{\J\per\square\m} \quad\Rightarrow\,
                     k=\qty{40}{\newton\per\meter}$
            \tcbitem (0,5pt)
            %--------------------------------------------------------
            \tcbitem $T_0=2\pi\sqrt{\frac{m}{k}}=2\pi\sqrt{\frac{0,25}{40}}
                     \simeq\qty{0.50}{\second}$
            \tcbitem (0,5pt)
        \end{tcbitemize}
    %----------------------------------------------------------------
    \tcbitem[raster multirow=2] 3-
    \tcbitem[raster multicolumn=2,raster multirow=2,blankest]
        \begin{tcbitemize}[raster rows=2,raster columns=2,
               raster height=\tcbtextheight,
               raster column 1/.style={width=\cwii},
               raster column 2/.style={width=\cwiii,fontupper=\bfseries}
            ]
            \tcbitem $E_m=E_{c}+E_{pe}=E_{pe,\max}=\frac12kX_m^2$
            \tcbitem (0,25pt)
            %--------------------------------------------------------
            \tcbitem $E_m=\frac12\times 40\times (\num{2e-2})^2=
                     \qty{8.0e-3}{\J}=\qty{8.0}{\mJ}$
            \tcbitem (0,25pt)
        \end{tcbitemize}
    %----------------------------------------------------------------
    \tcbitem[raster multirow=2] 4-
    \tcbitem[raster multicolumn=2,raster multirow=2,blankest]
        \begin{tcbitemize}[raster rows=2,raster columns=2,
               raster height=\tcbtextheight,
               raster column 1/.style={width=\cwii},
               raster column 2/.style={width=\cwiii,fontupper=\bfseries}
            ]
            \tcbitem $E_m=E_{c}+E_{pe}=\frac12mv_x^2+\frac12kx^2
                     \quad\Rightarrow\;
                     \frac12kX_m^2=
                     \frac12mv_x^2+\frac12k\left(\frac{X_m}{2}\right)^2$ \\
                     $\Rightarrow\,
                     v_x^2=\frac{k}{m}\left(X_m^2-\frac{X_m^2}{4}\right)=
                     \frac{3kX_m^2}{4m}
                     \quad\Rightarrow\;
                     v_{x}=\sqrt{\frac{3kX_m^2}{4m}}$
            \tcbitem (0,5pt)
            %--------------------------------------------------------
            \tcbitem le solide revient vers l'équilibre, la vitesse est
                     négative~: \\
                     $v_x=\sqrt{\frac{3\times40\times(\num{2e-2})^{2}}
                     {4\times0,25}}\simeq\qty{-0.22}{\v}$
            \tcbitem (0,25pt)
        \end{tcbitemize}
    %----------------------------------------------------------------
    \tcbitem[raster multirow=2] 5-
    \tcbitem[raster multicolumn=2,raster multirow=2,blankest]
        \begin{tcbitemize}[raster rows=2,raster columns=2,
               raster height=\tcbtextheight,
               raster column 1/.style={width=\cwii},
               raster column 2/.style={width=\cwiii,fontupper=\bfseries}
            ]
            \tcbitem $W(\vec{F})=\int_{0}^{\frac{X_m}{2}}-kx\,\mathrm{d}x=
                     \left[-\frac12kx^2\right]_{0}^{\frac{X_m}{2}}=
                     -\frac12k\left(\frac{X_m}{2}\right)^2$
            \tcbitem (0,5pt)
            %--------------------------------------------------------
            \tcbitem $W(\vec{F})=-\frac12\times 40\times (10^{-2})^2=
                     \qty{-2.0e-3}{\J}=\qty{-2.0}{\mJ}$
            \tcbitem (0,25pt)
        \end{tcbitemize}
\end{tcbitemize}


\end{document}
